% !TeX root = ../main.tex
\chapter{Durchführung und Auswertung des Experteninterviews}
\label{ch:experteninterview}

In diesem Kapitel wird die methodische Durchführung und Auswertung des qualitativen Experteninterviews 
dokumentiert. Die gewonnenen Erkenntnisse dienen innerhalb des gestaltungsorientierten 
Forschungsdesigns als empirische Brücke zum organisationalen und menschlichen Umfeld 
(\textit{Environment}). Sie bilden das Fundament zur Festlegung der manuellen Baseline sowie zur 
Definition der funktionalen Anforderungen an das Multi-Agenten-System.

\section{Zielsetzung und Erkenntnisinteresse}
\label{sec:zielsetzung}

Die Relevanz sowie die informationstechnische Kritikalität der Migration von LO-VC zu AVC wurden 
bereits in Kapitel~\ref{sec:lovc_context} dargelegt. Für die theoriegeleitete 
Konzeption eines \ac{MAS} ist die bloße Analyse theoretischer Syntaxunterschiede jedoch unzureichend. 
Spezifische, historisch gewachsene Problemstellungen in der Entwicklungspraxis weisen oft eine hohe 
Kontextabhängigkeit auf und lassen sich nicht rein analytisch aus Systemdokumentationen ableiten. 
Zudem mangelt es in der Literatur an quantitativen Daten bezüglich des tatsächlichen manuellen 
Migrationsaufwands.

Folglich ist die Durchführung eines Experteninterviews zur Beantwortung der Unter-Forschungsfragen 
zwingend erforderlich. Das primäre Erkenntnisinteresse konzentriert sich auf zwei Dimensionen:
\begin{itemize}
    \item \textbf{Quantifizierung der Baseline (FF1):} Erhebung valider zeitlicher Metriken und 
    prozessualer Schritte der manuellen Migration, um einen empirischen Vergleichsmaßstab für die 
    spätere Evaluation des Artefakts zu etablieren.
    \item \textbf{Anforderungsdefinition und Scoping (FF3):} Identifikation impliziter Praxisprobleme, 
    kognitiver Hürden und typischer Fehlerquellen im msg-Kontext, um daraus spezialisierte agentische 
    Rollen und funktionale Anforderungen für die Systemarchitektur abzuleiten.
\end{itemize}

\section{Methodisches Forschungsdesign}
\label{sec:forschungsdesign}

Um dem Anspruch forschungsmethodischer Strenge (\textit{Rigor}) gerecht zu werden, ist das 
Experteninterview als integraler Bestandteil der explorativen Fallstudie nach den Richtlinien von 
\textcite{runesonGuidelinesConductingReporting2009} gestaltet. Aufgrund der Komplexität der 
Code-Transformation, die maßgeblich durch undokumentiertes Erfahrungswissen (\textit{Tacit Knowledge}) 
geprägt ist, bedarf es zur fundierten Durchdringung der Problemstellung eines qualitativen 
Forschungsansatzes \parencite{seamanQualitativeMethodsEmpirical1999}. 

Als Erhebungsinstrument wird das semi-strukturierte Interview gewählt. Dieses Format kombiniert offene 
und spezifische Fragen in einem flexiblen Leitfaden. Dadurch wird sichergestellt, dass einerseits 
antizipierte Daten -- wie die zeitlichen Metriken für die Baseline -- systematisch erfasst werden, 
während andererseits ausreichend Raum für unvorhergesehene qualitative Erkenntnisse aus der 
Migrationspraxis bleibt \parencite{seamanQualitativeMethodsEmpirical1999}. Während in 
sozialwissenschaftlichen Disziplinen oft ein geringes Vorwissen der interviewenden Person gefordert 
wird, postulieren \textcite{hoveExperiencesConductingSemistructured2005} für den Bereich des Software 
Engineerings die Notwendigkeit eines tiefgreifenden technischen Vorverständnisses, um fachspezifische 
Nuancen während des Gesprächs korrekt zu erfassen und zu vertiefen.

Die Auswertung des Datenmaterials folgt einem systematischen Mixed-Methods-Ansatz, der Methoden der 
qualitativen und quantitativen Inhaltsanalyse kombiniert. Der Codierungsprozess gliedert sich in 
Anlehnung an \textcite{mayringQualitativeContentAnalysis2014} in zwei Phasen:
\begin{enumerate}
    \item \textbf{Deduktive Kategorienanwendung:} Vorab definierte Codes, die direkt aus den 
    Forschungsfragen abgeleitet wurden (z.\,B. manuelle Prozessschritte, Zeitaufwand), werden 
    theoriegeleitet auf das Textmaterial angewendet.
    \item \textbf{Induktive Kategorienbildung:} Ergänzende Codes für unvorhergesehene praxisnahe 
    Schmerzpunkte oder spezifische Syntaxprobleme werden iterativ direkt am empirischen Material 
    entwickelt.
\end{enumerate}

Im Sinne der in Kapitel~\ref{sec:methodisches_vorgehen} beschriebenen Datentriangulation bilden die 
Ergebnisse dieses Kapitels eine von drei empirischen Säulen. Sie werden im weiteren Verlauf der Arbeit 
mit den Erkenntnissen der Archivdokumentenanalyse sowie den Ergebnissen des praktischen Selbstversuchs 
synthetisiert, um die initiale, valide Architekturversion 1 (V1) des \ac{MAS} zu konstruieren. Um 
hierbei ein zentrales Qualitätsmerkmal qualitativer Fallstudien zu erfüllen, wird über den gesamten 
Auswertungsprozess -- von der Audioaufzeichnung über das codierte Transkript bis hin zu den daraus 
resultierenden Field Memos und funktionalen Anforderungen -- eine lückenlose Beweiskette 
(\textit{Chain of Evidence}) aufrechterhalten \parencite{runesonGuidelinesConductingReporting2009}. 
Dies garantiert die transparente und wissenschaftlich nachvollziehbare Ableitung jeder späteren 
agentischen Rolle aus den empirischen Rohdaten.

\section{Datenerhebung und Auswertungsprozess}
\label{sec:datenerhebung}

Das Experteninterview wurde am 17.03.2026 mit einem erfahrenen Product Owner und Softwareentwickler 
der \textit{msg systems ag} durchgeführt. Die Fachexpertise des Interviewpartners umfasst eine 
zwölfjährige Tätigkeit als Maschinenbauingenieur sowie eine mehr als sechsjährige Spezialisierung auf 
dem Gebiet der SAP-Variantenkonfiguration. Seit knapp vier Jahren liegt sein primärer 
Projektschwerpunkt auf der Betreuung von Kundenmigrationen von LO-VC nach AVC im msg-Kontext. Das 
Gespräch wies eine Dauer von rund 45 Minuten auf und wurde vom Autor dieser Arbeit auf Basis des 
vorab definierten Leitfadens moderiert.

Mit ausdrücklichem Einverständnis des Experten wurde das Gespräch digital aufgezeichnet und über die 
integrierte Funktionalität von Microsoft Teams automatisch transkribiert. Zur Sicherung der 
Datenqualität und Wahrung der wissenschaftlichen Integrität durchlief das Rohmaterial im Nachgang 
einen dreistufigen Aufbereitungsprozess:
\begin{enumerate}
    \item \textbf{Anonymisierung:} Manuelle Entfernung aller personenbezogenen und projektspezifischen 
    Identifikationsmerkmale zur Gewährleistung der Vertraulichkeit.
    \item \textbf{KI-gestützte Glättung:} Bereinigung des Transkripts von rein sprachlichen 
    Redundanzen (z.\,B. Füllwörter, Satzabbrüche) mittels generativer KI zur Optimierung der 
    Lesbarkeit.
    \item \textbf{Manuelle Verifikation:} Eine vollständige, zeilenweise Überprüfung des geglätteten 
    Textes durch den Autor, um sicherzustellen, dass durch die Bearbeitung keinerlei inhaltliche oder 
    technologische Verzerrungen eingeführt wurden.
\end{enumerate}

Das finale, geglättete Transkript (siehe Anhang~\ref{app:codierte_passagen}) bildete die 
Datengrundlage für das in Kapitel~\ref{sec:forschungsdesign} beschriebene Codierungsverfahren. 
Die deduktive Basis-Codeliste strukturierte das Material zunächst in die Hauptkategorien 
\textit{Prozessschritte der Migration}, \textit{Zeitaufwand und quantitative Metriken} sowie 
\textit{kognitive Hürden}. Durch den induktiven Codierdurchlauf konnten zusätzliche, spezifische 
Hürden -- wie etwa das fehleranfällige manuelle Parsing komplexer SAP-Prozeduren -- als kritische 
Systemanforderungen extrahiert werden. Die im Folgenden dargestellten Ergebnisse und Memos bilden die 
aggregierte Synthese dieses empirischen Prozesses.

\section{Ergebnisse der Befragung}
\label{sec:ergebnisse_interview}
Die nachfolgende Darstellung fasst die zentralen Erkenntnisse der qualitativen Inhaltsanalyse zusammen, welche durch die Synthese der Field Memos in thematische Schwerpunkte gegliedert wurden. Diese deskriptiven Ergebnisse bilden den empirischen Ist-Zustand des Migrationsprozesses ab und dienen als Grundlage für die spätere Anforderungsanalyse und die Definition der menschlichen Baseline.

\subsection{Kognitive Altlasten und historischer Code}
Die Auswertung der Interviewdaten verdeutlicht, dass historisch gewachsene LOVC-Modelle durch eine hohe Dichte an ungenutzten Altlasten charakterisiert sind. Merkmale, Prozeduren und Beziehungswissen verbleiben oft im System, ohne im aktuellen Modellierungskontext eine funktional Relevanz zu besitzen. Zudem führen über Jahre stetig anwachsende Werteräume zu einer erhöhten Komplexität. Für den Migrationsprozess bedeutet dies, dass eine unreflektierte 1:1-Konvertierung in den Advanced Variant Configurator (AVC) aufgrund dieser behindernden Altlasten in der Praxis unzureichend ist. Für den Entwickler resultiert daraus eine erhebliche kognitive Belastung, da eine zeitintensive manuelle Analyse zur Trennung von aktiver Logik und „totem Code“ zwingend erforderlich ist, bevor die eigentliche syntaktische Transformation beginnen kann.

\subsection{Paradigmenwechsel und syntaktisches Umdenken}
Über die rein syntaktische Ebene hinaus zeigt die Befragung einen fundamentalen Paradigmenwechsel in der Modellierungstechnik auf. Dieser wird im industriellen Kontext unter anderem durch moderne SAP-Architekturvorgaben wie die „Clean Core“-Strategie flankiert, wodurch proprietäre Eigenentwicklungen wie Z-Functions im AVC nicht mehr unterstützt werden. Da sich das Systemverhalten im AVC selbst bei identischer Syntax massiv ändern kann – etwa durch den Zwang zu klassifizierten Merkmalen oder die globale Wirkung von Constraints – führt eine bloße Übernahme alter Logiken häufig zu abweichenden Konfigurationsergebnissen oder Performance-Einbußen. Der Migrationsprozess erfordert daher ein aktives Umdenken: Entwickler müssen die ursprüngliche Intention des Legacy-Codes in Gänze verstehen und diese semantisch mit den neuen Bordmitteln des AVC abbilden.

\subsection{Fragmentierung manueller Prozessschritte}
Trotz der Existenz von Standard-Tools zur Massenanlage von Objekten bleibt der Migrationsprozess in weiten Teilen fragmentiert und durch manuelle Arbeitsschritte geprägt. Während Automatisierungen die Anlage von Basisdaten übernehmen können, verbleibt die komplexe intellektuelle Durchdringung der Logik beim Entwickler. Dieser muss Vorbedingungen manuell auswerten, in tabellenbasierte Strukturen überführen und abschließend auf Konsistenz prüfen. Dieser Workflow erzwingt häufig Medienbrüche und degradiert den Experten zu einer „manuellen Übersetzungsmaschine“, was die Fehleranfälligkeit erhöht und den Prozess verlangsamt.

\subsection{Limitierungen bestehender Standard-Tools}
Die Auswertung zeigt, dass bestehende Standardlösungen wie das „SAP Migration Cockpit“ primär die rein technischen Datenübernahmen (wie die Massenanlage von Objekten) abdecken. Bei der komplexen Logik-Transformation entstehen jedoch funktionale Lücken. Der Experte äußert den Bedarf nach einer prozessualen Unterstützung auf einer „beratenden Flugebene“, welche Zusammenhänge aufzeigt und Optimierungspotenziale identifiziert. Es mangelt im aktuellen Ist-Zustand an Werkzeugen zur automatisierten Kommentierung von historischem Code, zur tiefgehenden Analyse von Z-Functions sowie an proaktiven Ratschlägen zur Modellstrukturierung. Eine rein statische Code-Konvertierung erweist sich für den Gesamtprozess somit als unzureichend.

\subsection{Manueller Aufwand und Fehleranfälligkeit im Testprozess}
Ein signifikanter Flaschenhals wurde im Bereich der Qualitätssicherung identifiziert. Das Testen migrierter Modelle erfolgt gegenwärtig weitgehend manuell, da bestehende Testtools oft nicht in der Lage sind, logische Veränderungen zwischen Alt- und Neusystem korrekt zu interpretieren. Dies führt zu einem massiven Aufwand beim händischen Abgleich von Testaufträgen und birgt ein entsprechendes Potenzial für menschliche Flüchtigkeitsfehler in der abschließenden Validierungsphase.

\subsection{Zeitliche Metriken und Identifikation von Flaschenhälsen}
Zur Definition einer menschlichen Baseline wurden im Rahmen der Befragung konkrete zeitliche Metriken erhoben. Die Migration eines mittleren Modells beansprucht demnach etwa drei bis vier Personentage, wobei Extremfälle bei komplexen Strukturen bis zu mehrere Monate beanspruchen können. Als primärer zeitlicher Flaschenhals wurde die Transformation der Vorbedingungswelt identifiziert, die etwa zwei Drittel (66 \%) der gesamten Arbeitszeit beansprucht. Weitere 40 \% des Aufwands entfallen auf die vorbereitende Analyse und den Aufbau unterstützender Strukturen wie Excel-Tabellen. Diese quantifizierten Werte bilden die Vergleichsvariablen für die spätere Evaluierung.

Die hier dargelegten empirischen Befunde skizzieren die aktuellen prozessualen und kognitiven Hürden der Migration. Sie bilden die konzeptionelle Ausgangslage, um im folgenden Unterkapitel konkrete Implikationen und Anforderungen an die zu entwickelnde Systemarchitektur abzuleiten.

\section{Implikationen für die Systemarchitektur (Ideales Zielbild)}
\label{sec:implikationen}
Basierend auf den empirischen Befunden der qualitativen Inhaltsanalyse lassen sich spezifische Anforderungen und konzeptionelle Leitplanken für die Gestaltung der Systemarchitektur ableiten. In diesem Kapitel werden ausschließlich jene Implikationen dargelegt, die direkt aus der Expertenbefragung hervorgehen. Sie repräsentieren den maximalen \textit{Business Need} der industriellen Praxis und skizzieren das theoretische Idealbild des Systems.

\subsection{Funktionale Anforderungen aus der Expertenperspektive}
Aus den identifizierten prozessualen Hürden ergeben sich aus Sicht der Praxis folgende zentrale funktionale Wünsche an das Systemdesign:
\begin{itemize}
    \item \textbf{Semantische Analyse und Code-Bereinigung:} Da historische Altlasten die Migration behindern, sollte die Architektur eine Komponente zur Identifikation von ungenutzten Merkmalen und Prozeduren enthalten. Ziel ist eine Reduktion der kognitiven Last durch eine automatisierte Vorab-Bereinigung des Legacy-Codes.
    \item \textbf{Automatisierte Logik-Transformation:} Angesichts des zeitlichen Flaschenhalses bei der Erstellung von Tabellen und Constraints muss der Fokus auf der automatisierten Generierung dieser Strukturen liegen. Das System sollte idealerweise in der Lage sein, die logische Restrukturierung von Vorbedingungen zu übernehmen, um den manuellen Aufwand drastisch zu minimieren.
    \item \textbf{Konsistenzprüfung im Testprozess:} Um die Fehleranfälligkeit bei der Validierung zu senken, wünscht sich die Praxis eine Funktion zum Mapping alter Testkonfigurationen auf die Zielmodelle, die logische Abweichungen zwischen den Systemwelten erkennt und den manuellen Testaufwand reduziert.
\end{itemize}

\subsection{Architektonische Paradigmen und Beratungsansatz}
Der identifizierte Paradigmenwechsel in der Modellierungstechnik bedingt zudem spezifische architektonische Eigenschaften:
\begin{itemize}
    \item \textbf{Interaktive Assistenz (Consulting-Ansatz):} Das System sollte nicht als rein deterministischer Konverter agieren. Gefordert ist eine interaktive Ebene, die Design-Entscheidungen begründet und dem Nutzer beratend zur Seite steht, um den Kontext der Transformation transparent zu machen.
    \item \textbf{Einhaltung von AVC-Modellierungsrichtlinien:} Die Architektur sollte standardkonformen AVC-Zielcode erzeugen. Dabei gilt es, grundlegende Performance-Best-Practices der SAP -- wie den Verzicht auf obsolete Z-Funktionen zugunsten nativer Standard-Funktionen -- direkt im Transformationsschritt zu berücksichtigen, um eine fehlerfreie Lauffähigkeit in der neuen Konfigurations-Engine sicherzustellen.
\end{itemize}

\subsection{Hypothese der Multi-Agenten-Rollen}
Die Heterogenität der Anforderungen – von der Analyse komplexer Z-Functions bis hin zur Validierung von Test-Sets – legt eine Verteilung der Aufgaben auf spezialisierte Rollen nahe. Ein monolithischer Ansatz erscheint unzureichend, um die benötigte Tiefe in den Bereichen Analyse, Beratung, Synthese und Test gleichzeitig abzudecken. Daraus ergibt sich die konzeptionelle Hypothese einer Multi-Agenten-Architektur, in der idealerweise spezialisierte Agenten (wie z.\,B. Analyst, Syntaktiker, Tester und Consultant) kollaborativ zusammenarbeiten. 

\subsection{Quantitative Zielvorgaben für die Systemvalidierung}
Die erhobenen Zeitmetriken definieren die harten Erfolgskriterien für die spätere Evaluierung der Architektur. Das System muss gegen die menschliche Baseline von drei bis vier Tagen pro Modell validiert werden. Insbesondere soll nachgewiesen werden, dass der primäre Flaschenhals der Vorbedingungs-Transformation (66\,\% des Gesamtaufwands) sowie die vorbereitende Analysephase (40\,\% des Aufwands) durch die automatisierte Unterstützung messbar reduziert werden können.

\subsection{Überleitung zum methodischen Realitätsabgleich}
\label{sec:ueberleitung_scoping}
Die hier dargelegten interviewbasierten Implikationen definieren das angestrebte Idealbild der Praxis. Im Sinne des methodischen \textit{Design Science Research} Vorgehens (vgl. Kapitel~\ref{sec:methodisches_vorgehen}) darf dieses vollumfängliche Anforderungsprofil jedoch nicht unreflektiert in den Systementwurf übernommen werden. Vielmehr bedarf es nun eines technologischen Realitätsabgleichs. Im darauffolgenden Kapitel (Selbstversuch) werden diese idealisierten Forderungen – insbesondere hochkomplexe Wünsche wie die automatisierte Test-Migration oder weitreichende Consulting-Funktionen – auf Code-Ebene auf ihre tatsächliche softwaretechnische Machbarkeit geprüft. Der Selbstversuch fungiert somit als methodischer Filter, um valide zu entscheiden, welche Anforderungen in die finale Agenten-Architektur (Version 1) aufgenommen werden und welche Aspekte aufgrund unverhältnismäßiger technologischer Komplexität bewusst aus dem \textit{Scope} dieser Masterarbeit exkludiert werden müssen.